% =====================================================================
%  methods-v3 / proposed-solution-blue-v3.tex
%  Manuscript-ready DRAFT of the Proposed Solution section (v3).
%  Deliverable for MICROENCODER_HARDENING_BRIEF.md Sec. 21.
%  Prepared 2026-10-04 (gnhf microencoder-hardening sprint, iteration 20).
%
%  WHAT THIS IS
%  ------------
%  Replacement for methods-v2/proposed-solution-blue.tex (left unchanged).
%  It is built from the v3 decisions in 01 (model-level utility), 02 (split
%  equivalence, Branches X/L/E), 03 (KD hardening, C* control), 08 (utility
%  definition), 09 (cost policy C-1) and 12 (Conservative wording), and cites
%  ONLY keys from method-references-v3.bib, within the "May support" limits
%  recorded in 17-verified-references-v3.md Sec. 2 and the usage rules of
%  17 Sec. 5. Paragraph tags "% [Pn]" are the anchors used by
%  16-citation-audit.md.
%
%  BLUE CONVENTION
%  ---------------
%  Everything in this section is new manuscript text, so the whole section
%  body is typeset in blue (\color{blue} at the start of the section and
%  inside each float, because colour does not propagate into floats).
%
%  HOW TO MERGE INTO THE IEEEtran MANUSCRIPT (do NOT edit the original here)
%  ----------------------------------------------------------------------
%    1. Add \usepackage{xcolor} to the manuscript preamble.
%    2. Copy the text between BEGIN MERGE and END MERGE into the empty
%       \section{Proposed Solution} of the manuscript.
%    3. Append method-references-v3.bib entries to the manuscript .bib.
%    4. Equation numbers continue after the manuscript's Eqs. (1)-(3);
%       f_theta, M_k, lambda and C keep the manuscript's meaning.
%
%  COMPILATION STATUS
%  ------------------
%  Standalone article-class wrapper so the section can be proofed alone.
%  No LaTeX toolchain (pdflatex/latexmk/bibtex) exists in the sprint
%  environment, so this file has NOT been compiled. Only a static check
%  (balanced braces/environments, every \cite key present in the .bib) was
%  run. It must be compiled at merge time.
%
%  WORDING RULES APPLIED (brief Sec. 10, 21, 30)
%  ---------------------------------------------
%  No "first" claim; no "without retraining" claim; no claim that network
%  changes force retraining; split quality is never asserted equal, only
%  tested (E0); KD is investigated, not assumed to help; continual learning
%  is out of scope; no results are reported (none exist yet).
% =====================================================================
\documentclass[10pt]{article}
\usepackage{amsmath,amssymb,amsfonts}
\usepackage{algorithm}
\usepackage{algorithmic}
\usepackage{xcolor}
\usepackage{cite}
\usepackage{url}

\begin{document}

% ======================== BEGIN MERGE ========================
\section{Proposed Solution}
\label{sec:proposed}
\color{blue}

% [P1] Design overview
\subsection{Design Overview}
We propose a router for edge--cloud LLM serving in which a query can be
answered by a small language model (SLM) on the edge device, by split
execution of a larger model across the edge device and a server, or by
remote execution on the server. The router separates two questions that
have different sources of evidence. \emph{Which model is likely to answer
this query well} is a semantic question, answered by a compact text encoder
trained offline on observed model outcomes. \emph{What it costs to run that
model now, and where} is a systems question, answered by a cost model that is
filled with profiled compute times and live network and load measurements at
decision time. A runtime controller combines the two to choose an execution
action. The design builds on learned query routers that predict model quality
from the query text~\cite{ong2025routellm,ding2024hybridllm,chen2024routerdc,ding2025bestroute},
on LLM routing at the network edge that trades predicted quality against
latency or communication cost~\cite{bao2025qualitylatency,chen2025netgptrouting,xue2026cr2}
or tunes routing thresholds for latency across device, edge and cloud tiers~\cite{qiao2026consroute},
and on split computing, which selects a partition point of a neural network
under bandwidth and server-load conditions~\cite{kang2017neurosurgeon,li2020edgent,laskaridis2020spinn,matsubara2022splitsurvey}.
The remainder of this section defines the execution actions
(Section~\ref{sec:actions}), the utility the router predicts
(Section~\ref{sec:utility}), the encoder (Section~\ref{sec:encoder}), its
training (Section~\ref{sec:kd}), the runtime cost and decision rule
(Section~\ref{sec:cost}) and the split execution path
(Section~\ref{sec:split}).

% [P2] Actions and models
\subsection{Execution Actions and the Action-to-Model Map}
\label{sec:actions}
Let the action set be
\begin{equation}
\mathcal{A} = \{\textsc{Edge},\ \textsc{Split}(d_1),\ \ldots,\
\textsc{Split}(d_m),\ \textsc{Remote}\},
\label{eq:actions}
\end{equation}
where $\textsc{Split}(d)$ runs the first $d$ decoder layers of a model on
the edge device and the remaining layers on the server. We distinguish an
\emph{action} from the \emph{model} it executes. A model specification $M$
fixes the checkpoint and revision, the weight quantization, the compute data
type, the prompt template and the decoding configuration; the same
checkpoint at two numerical configurations is two specifications. Each action
$a$ executes exactly one specification $M(a)$. We write
$\mathcal{M} = \{M(a) : a \in \mathcal{A}\}$ for the candidate models, with
$K = |\mathcal{M}|$, and $\mathcal{A}(M) = \{a : M(a) = M\}$ for the
realizations of model $M$. All split depths of the split-capable model $M_S$
are realizations of $M_S$; if the remote action runs the same specification,
it is the realization in which all layers run on the server. The candidate
models $\mathcal{M}$ play the role of the pool $\mathcal{M}$ in the problem
formulation, so the scoring function $f_\theta(q, M_k)$ of
Eq.~(3) remains a function of a query and a \emph{model}.

% [P3] Utility definition
\subsection{Utility: Expected Task Score of a Model on a Query}
\label{sec:utility}
For a dataset or task family $D$ we fix, before training, a score function
$S_D(y, q) \in [0,1]$ computed by the benchmark's own grader (exact match,
unit tests, or a pinned judge with a fixed prompt and version) from the output
$y$ and the reference answer of $q$. The utility of model $M$ on query $q$ is
the expected score
\begin{equation}
U(M \mid q) = \mathbb{E}_{\xi}\big[S_D\big(y_M(q;\xi),\, q\big)\big] \in [0,1],
\label{eq:utility}
\end{equation}
where $\xi$ collects the randomness of a run: sampling during decoding,
run-to-run numerical variation, and grader stochasticity. For a binary metric
$U(M\mid q)$ is the probability that $M$ answers $q$ correctly; for a graded
metric it is the expected score. On SimpleQA, for example, a response scores
1 only if the benchmark's grader marks it correct, and both incorrect answers
and abstentions score 0~\cite{wei2024simpleqa}. We do not define a single
notion of ``success'' across heterogeneous benchmarks: each dataset keeps its
own $S_D$, and results are reported per dataset and macro-averaged.
Utility depends on the query and the model only; it does not depend on
network state, cost, or the trade-off weight $\lambda$.

% [P4] Labels
Training labels are obtained by running every model specification in
$\mathcal{M}$ on the training queries and scoring the outputs with $S_D$. A
single run gives a noisy but unbiased estimate of $U(M\mid q)$; where the
budget allows, we average $k$ sampled runs, following the use of empirical
soft labels from repeated sampling in Hybrid LLM~\cite{ding2024hybridllm}.
Labels for the edge SLM and the split-capable model are produced by new
inference runs through the benchmark's evaluation pipeline, because
pre-computed benchmark results do not cover these model
specifications.

% [P5] Encoder
\subsection{The Micro-Encoder Scoring Function $f_\theta$}
\label{sec:encoder}
The router is a compact bidirectional encoder $E_\theta$ followed by one
sigmoid head per candidate model:
\begin{equation}
f_\theta(q, M_k) := \hat U_\theta(M_k \mid q)
= \sigma\!\big(w_k^{\top} h_\theta(q) + b_k\big),\qquad
h_\theta(q) = \mathrm{pool}\big(E_\theta(q)\big),
\label{eq:ftheta}
\end{equation}
for $k = 1,\ldots,K$, with mean pooling over token states. The heads are
independent sigmoids rather than a softmax over models, because the events
``model $M_k$ answers $q$ correctly'' are not mutually exclusive: several
models may all answer a query correctly, or all fail. Per-model sigmoid
outputs of this kind are also used by the prefill-activation router of
Varshney et al.~\cite{varshney2026prefill} and by the CR$^2$ router of
Xue et al.~\cite{xue2026cr2} (both preprints), and BEST-Route trains one
sigmoid head per (model, number-of-samples) action~\cite{ding2025bestroute}.

% [P6] Backbone
The encoder is a replaceable component, and its choice is an engineering
decision rather than a contribution. Our first implementation candidate is
ModernBERT-base, an encoder with 149M parameters, 22 layers and a native
context of 8{,}192 tokens~\cite{warner2025modernbert}; ModernBERT-base was
also used as the encoder of the HyDRA router (preprint)~\cite{garg2026hydra}.
The compared alternatives are smaller encoders of the Ettin
suite~\cite{weller2026ettin}, which replicate the ModernBERT recipe at sizes
down to 17M parameters, DeBERTa-v3-small~\cite{he2023debertav3}, the
encoder family used by the Hybrid LLM and BEST-Route
routers~\cite{ding2024hybridllm,ding2025bestroute}, and a MiniLM
encoder~\cite{wang2020minilm}. Li et al.\ report only a limited effect of the
embedding model in LLMRouterBench, for three frozen embedding models
combined with embedding-based routers~\cite{li2026llmrouterbench}; that
result does not cover fine-tuned encoders, so we treat the backbone as an
experimental factor. Each candidate is trained with the same non-distilled
protocol and tuning budget, and the deployed backbone is the
lowest-latency candidate whose validation routing utility lies within a
tolerance, fixed in advance, of the best candidate.

% [P7] Pool changes
Because the heads in Eq.~\eqref{eq:ftheta} are indexed by model, adding a
new split depth of an existing model adds no head: it requires the
split-equivalence test of Section~\ref{sec:split} at that depth and a cost
profile, and, if the test passes, no new utility labels. Adding a new
\emph{model} requires labelled outcomes for that model, training at least a
new head on the encoder features, and recalibration. Routing to models never
seen in training would require model descriptors as router inputs, as in
IRT-Router~\cite{song2025irtrouter}, the learned model embeddings of
RouterDC~\cite{chen2024routerdc}, or the capability profiles of HyDRA
(preprint)~\cite{garg2026hydra}; our design does not provide this.

% [P8] Teacher
\subsection{Teacher--Student Training}
\label{sec:kd}
We investigate whether distilling a larger router improves the compact
encoder over a well-tuned student trained without distillation. The
teacher is a larger encoder from the same tokenizer family (ModernBERT-large
or an Ettin encoder of similar size) with the same per-model heads, trained
on the same labels. Since such a teacher sees no information beyond the
student's labels, any benefit must come from statistical effects, such as
reducing the variance of noisy single-run labels~\cite{menon2021statistical}
or acting as a learned, query-dependent form of label
smoothing~\cite{yuan2020revisiting}, and gains are not guaranteed: students
often do not match their teachers~\cite{stanton2021kd}, and larger teachers
are not always better teachers~\cite{cho2019efficacy}. Teacher size is
therefore an experimental factor. A teacher that reads privileged
training-time inputs, such as the edge SLM's prefill activations, is a
generalized-distillation variant~\cite{lopezpaz2016unifying} that we treat as
a secondary condition. Distilling a router into a smaller on-device router has
been used before: CR$^2$ (preprint) distils an on-device gate from a teacher
with privileged information~\cite{xue2026cr2}.

% [P9] Cross-fitting and teacher calibration
Teacher targets on training queries are produced \emph{out of fold}: the
training set is split into $F$ folds, $F$ teachers are each trained on
$F-1$ folds, and every query receives its target from the teacher that did
not see it. Without this step a high-capacity teacher can fit its training
queries almost exactly, and its targets collapse onto the hard labels. Each
teacher head is then calibrated with Platt scaling~\cite{guo2017calibration}
on a teacher-calibration split that is disjoint from the student's
calibration and test splits, so that the soft targets
$\tilde t(q, M) \in [0,1]$ estimate the expected score
of Eq.~\eqref{eq:utility}.

% [P10] Student loss
The student minimizes a weighted sum of a hard-label term and a soft-target
term, the form introduced for knowledge distillation by Hinton
et al.~\cite{hinton2015distilling}:
\begin{equation}
\mathcal{L}_S(\theta) = \sum_{i}\sum_{M\in\mathcal{M}} w_{D(i)}\Big[
\alpha\,\ell\big(\hat U_\theta(M\mid q_i),\, y_i(M)\big)
+ (1-\alpha)\,\ell\big(\hat U_\theta(M\mid q_i),\, \tilde t(q_i, M)\big)\Big],
\label{eq:kdloss}
\end{equation}
where $\ell(p,t) = -[t\log p + (1-t)\log(1-p)]$ is the binary cross-entropy
with a target $t\in[0,1]$, $y_i(M)$ is the observed score, and $w_D$ is a
per-dataset weight. Because $\ell$ is linear in $t$, Eq.~\eqref{eq:kdloss}
equals the cross-entropy against the mixed target
$\alpha y_i(M) + (1-\alpha)\tilde t(q_i,M)$, and since the expected
cross-entropy is minimized at $p = \mathbb{E}[t]$, the student is trained
towards the conditional mean that the mixed target encodes. This makes the
comparison with non-distilled training exact: label
smoothing~\cite{muller2019labelsmoothing} mixes in a constant, smoothing
towards a per-dataset base rate mixes in a per-model prior, and distillation
mixes in a query-dependent teacher output.

% [P11] Controls and terminology
The non-distilled control is the same student trained on the same labels,
with label smoothing and base-rate smoothing in its hyperparameter search,
weight decay and early stopping, given the same number of tuning trials as
the distilled student. A self-distillation control, in which the teacher has
the student's size~\cite{furlanello2018bornagain}, separates the effect of a
larger teacher from that of any teacher. Every student, distilled or not,
is calibrated after training with per-model Platt or temperature
scaling~\cite{guo2017calibration} on the same calibration split, so any
calibration benefit of distillation must exceed that of post-hoc calibration.
If distillation does not improve routing utility over this control, the
non-distilled student is deployed. We use the term distillation only for
training on a teacher router's soft outputs. RouteLLM's augmentation of human
preference data with labels from GPT-4 acting as a judge~\cite{ong2025routellm,zheng2023judging}
is preference supervision with hard labels; training on teacher-generated
hard targets is a recognized but distinct sequence-level
variant~\cite{kim2016sequence}.

% [P12] Cost and state
\subsection{Runtime Cost and Decision Rule}
\label{sec:cost}
At decision time the edge node measures a system state $s$: round-trip time
to the server, available uplink and downlink bandwidth, the server's queue
length, the server's cold or warm status, and edge utilization and free
memory. The cost of an action is its predicted end-to-end latency
$\hat L(a\mid q,s)$, assembled from the time-to-first-token and
time-per-output-token decomposition
$T_{\mathrm{total}} = \mathrm{TTFT} + (N-1)\,\mathrm{TPOT}$, where the compute
terms come from offline profiles of each action on the deployment hardware
and the communication terms are computed from the measured state. For
$\textsc{Split}(d)$ the communication terms include the transfer of the
activation tensor at the cut, $B/\mathrm{BW}_{\mathrm{up}}$, and the
round trips needed per decoding step. The expected output length $N$ is the
per-model median on a profiling split. Transferred bytes, energy, monetary
cost and memory are measured and reported separately rather than summed with
latency.

% [P13] Decision rule
The controller chooses
\begin{equation}
R(q,s) = \operatorname*{arg\,max}_{a \in \mathcal{F}(q,s)}
\Big[\hat U_\theta\big(M(a)\mid q\big) - \hat\Delta(a,q)
- \lambda\, \hat L(a\mid q,s)\Big],
\label{eq:route}
\end{equation}
where $\hat\Delta(a,q)$ is a quality correction for split execution that is
zero unless the equivalence test of Section~\ref{sec:split} fails, and the
feasible set
\begin{equation}
\mathcal{F}(q,s) = \{a\in\mathcal{A}:\ \hat L(a\mid q,s)\le L_{\mathrm{SLA}},\
\mathrm{mem}_{\mathrm{edge}}(a)\le \mathrm{mem}_{\mathrm{avail}}(s),\
a \text{ permitted by the data policy}\}
\label{eq:feasible}
\end{equation}
applies a per-query latency limit, the edge memory limit and any rule that
forbids raw text from leaving the device. With $\hat\Delta \equiv 0$ and
$\mathcal{F}=\mathcal{A}$, Eq.~\eqref{eq:route} instantiates the routing
objective of the problem formulation with
$f_\theta(q,M_k) = \hat U_\theta(M_k\mid q)$ and $C = \hat L$.
Additive quality-minus-cost rules of this form appear in prior
routers~\cite{song2025irtrouter,varshney2026prefill}, including routers
at the wireless edge that subtract a latency or communication
penalty~\cite{bao2025qualitylatency,xue2026cr2}. If $\mathcal{F}(q,s)$ is
empty, the controller runs \textsc{Edge} when it fits in memory and logs a
latency-limit violation.

% [P14] lambda
Since the objective is an expectation over queries of a term that depends
only on each query and its action, choosing the action per query as in
Eq.~\eqref{eq:route} maximizes the plug-in objective for a fixed $\lambda$.
Over randomized policies, $\lambda$ can also be read as the Lagrange
multiplier of an average-latency budget; the per-query limit
$L_{\mathrm{SLA}}$ is enforced as a feasibility constraint and is not
dualized. The weight $\lambda$, in units of normalized score per second, is
swept to trace quality--latency frontiers. For headline results it is chosen
on the validation split to meet a stated mean-latency or remote-call budget
and then frozen. We do not claim constrained routing as a contribution.

% [P15] Why s stays out of theta
Network state enters the decision only through the measured cost term, so the
router reacts to changes in bandwidth and round-trip time at the next query
without updating the semantic predictor. A router that takes $s$ as an input
could learn this dependence from data, but it would need training data
spanning the network conditions it must handle. Keeping the two parts separate
lets one semantic predictor be reused across deployments, lets the cost
model be re-profiled without new utility labels, and makes each routing
decision attributable to $\hat U_\theta$, $\hat\Delta$ and $\lambda\hat L$.

% [P16] Split execution and equivalence
\subsection{Split-LLM Execution and Split Equivalence}
\label{sec:split}
In $\textsc{Split}(d)$ the edge device runs the embedding layer and the
first $d$ decoder layers of $M_S$ and sends the hidden states at the cut to
the server, which runs the remaining layers and returns tokens. For a
standard dense decoder the tensor crossing the cut has
$n \times d_{\mathrm{model}}$ elements for $n$ prompt tokens and hidden size
$d_{\mathrm{model}}$, so the prefill transfer size does not depend on the
depth $d$; depth mainly moves compute between the edge and the server.
Prior systems select split points of neural networks under measured
bandwidth and load~\cite{kang2017neurosurgeon,li2020edgent,laskaridis2020spinn},
and recent work learns split points for LLMs under changing channel or
bandwidth conditions for a single model~\cite{chen2025adaptivesplit,younesi2025splitwise}.

% [P17] When depths share utility
Whether split depths may share the utility of their model is a hypothesis
that we test, not an assumption. Exact tensor transfer does not by itself
guarantee bit-identical outputs, because floating-point results depend on
batch size, hardware and kernel choice, which can change greedy LLM
outputs even without splitting~\cite{yuan2025nondeterminism}; a technical
report attributes such run-to-run differences to kernels that are not
batch-invariant~\cite{he2025nondeterminism}. Quantizing or compressing the
transmitted activations~\cite{laskaridis2020spinn} or adding noise to them for
privacy~\cite{mai2024splitdenoise} makes split execution approximate by
design, and early exit changes inference semantics, with outputs that are
consistent with the full model only within a
tolerance~\cite{schuster2022calm}. Before training the router we therefore
run a split-equivalence test (E0) at every deployed depth and numerical
configuration: per-query task scores of $\textsc{Split}(d)$ are compared with
those of unsplit execution of the same specification, against the unsplit
model's own run-to-run variation. Three branches follow.
\emph{Exact:} if E0 passes, every depth of $M_S$ shares
$\hat U_\theta(M_S\mid q)$ and $\hat\Delta \equiv 0$, so
\begin{equation}
\max_{a\in\mathcal{A}}\big[\hat U_\theta(M(a)\mid q) - \lambda \hat L(a\mid q,s)\big]
= \max_{M\in\mathcal{M}}\Big[\hat U_\theta(M\mid q)
- \lambda \min_{a\in\mathcal{A}(M)} \hat L(a\mid q,s)\Big],
\label{eq:factor}
\end{equation}
that is, the router chooses a model and the cost model alone chooses its
depth. \emph{Lossy:} if E0 shows a quality loss at some depths, we subtract a
per-depth correction $\hat\Delta(d)$ estimated on a held-out set, and move to
a query-dependent correction $\hat\Delta(d,q)$, or to separate heads per
depth, only if the loss varies across queries. \emph{Semantics-changing:} an
execution path that alters the computation, such as early exit, is treated as
a different model specification with its own head and labels.

% [P18] Split cost caveat
Under a pipeline split in which the server samples tokens, every decoding
step needs a round trip between the edge and the server. When the pool also
contains remote execution of the same specification and the server is faster
than the edge device, a latency-only cost will then rarely select
$\textsc{Split}(d)$; split execution is chosen when the data policy forbids
sending raw text or a server-compute budget applies. We report how often
each action is selected under each condition.

% [P19] Algorithms
\subsection{Training and Inference Procedures}
Algorithm~\ref{alg:train} summarizes training and Algorithm~\ref{alg:infer}
summarizes routing of one query.

\begin{algorithm}[t]
\color{blue}
\caption{Router training (offline)}
\label{alg:train}
\begin{algorithmic}[1]
\REQUIRE training, teacher-calibration, calibration and validation splits;
model specifications $\mathcal{M}$; split depths $D$; folds $F$
\STATE Run E0 for every $d\in D$; set the branch (exact, lossy or
semantics-changing) and, if lossy, estimate $\hat\Delta$
\STATE For each query $q_i$ and each $M\in\mathcal{M}$, run $M$ and record
$y_i(M) = S_D(y_M(q_i), q_i)$ (or a $k$-run mean)
\FOR{$f = 1$ \TO $F$}
\STATE Train teacher $T_f$ on all folds except $f$; Platt-calibrate each head
on the teacher-calibration split
\STATE For queries in fold $f$, set $\tilde t(q_i, M) \leftarrow T_f(q_i, M)$
\ENDFOR
\STATE Train the student with Eq.~\eqref{eq:kdloss}, tuning $\alpha$ on the
validation split; train the non-distilled control with the same budget
\STATE Calibrate every student head with Platt or temperature scaling on the
calibration split
\STATE Keep the distilled student only if it beats the control on
validation routing utility; otherwise keep the control
\STATE Profile compute times of every $a\in\mathcal{A}$ on the deployment
hardware
\end{algorithmic}
\end{algorithm}

\begin{algorithm}[t]
\color{blue}
\caption{Routing one query (online)}
\label{alg:infer}
\begin{algorithmic}[1]
\REQUIRE query $q$; calibrated student; cost profiles; $\lambda$;
$L_{\mathrm{SLA}}$; data policy
\STATE Compute $\hat U_\theta(M\mid q)$ for all $M\in\mathcal{M}$ with one
encoder pass
\STATE Read the current state $s$ (RTT, bandwidth, server queue, cold or warm
status, edge memory)
\STATE For each $a\in\mathcal{A}$, compute $\hat L(a\mid q,s)$ and test
membership in $\mathcal{F}(q,s)$
\IF{$\mathcal{F}(q,s) = \emptyset$}
\STATE Run \textsc{Edge} if it fits in memory; log the violation
\ELSE
\STATE Select $a^\star$ by Eq.~\eqref{eq:route}; break ties by lower server
compute, then fewer transferred bytes
\ENDIF
\STATE Execute $a^\star$; log $\hat U_\theta$, $\hat L$, $s$, the measured
TTFT, TPOT, bytes and the output
\end{algorithmic}
\end{algorithm}

% [P20] Assumptions and limitations
\subsection{Assumptions and Limitations}
The method assumes a fixed pool of model specifications chosen before
training; changes to the pool after deployment are outside the scope of this
paper. Sharing one utility estimate across split depths is valid only where
E0 passes, and the lossy and semantics-changing branches require
depth-specific labels. Distillation is evaluated against a tuned and
calibrated non-distilled control and is not required by the rest of the
design. The latency model is analytic and its accuracy is measured
separately; decisions are only as good as the calibration of
$\hat U_\theta$ and the accuracy of $\hat L$. Several of the closest related
routers (HyDRA, CR$^2$, the prefill-activation router and the NetGPT
router) are preprints, and comparisons with them are limited to their
published descriptions. Our search for prior work combining query-aware
model selection, split execution and live network cost was bounded, and we
do not claim that the combination is new.

\color{black}
% ========================= END MERGE =========================

\bibliographystyle{IEEEtran}
\bibliography{method-references-v3}

\end{document}
